% Zone „Das kennst du schon“ – aus bank/einheiten/zone.jsonl (zusammenbau v0.1)
\einheitenkopf[zone][Kennst du schon]{Das kennst du schon}
\setcounter{aufgabe}{0}

%% TODO Titel: Ich-kann-Satz fehlt in der Bank (Platzhalter: Kettenname) – Nr. 1
%% TODO Anweisung: ein Satz über der ersten Teilaufgabe fehlt in der Bank – Nr. 1
\begin{aufgabe}{Multiplizieren und Dividieren mit zehn, hundert und tausend als Kommaverschiebung, auch bei Dezimalzahlen}
\begin{teile}
\teil $4{,}2$ mal hundert? \leerfeld
\teil $0{,}36$ mal tausend? \leerfeld
\end{teile}
\end{aufgabe}

%% TODO Titel: Ich-kann-Satz fehlt in der Bank (Platzhalter: Kettenname) – Nr. 2
%% TODO Anweisung: ein Satz über der ersten Teilaufgabe fehlt in der Bank – Nr. 2
\begin{aufgabe}{Stellenwerttafel lesen und schreiben, Nullen an der richtigen Stelle ergänzen (null Komma null sechs Meter; eintausendfünfhundert Milliliter)}
\begin{teile}
\teil $3$ Einer, $4$ Zehntel, $7$ Hundertstel – welche Zahl? \leerfeld
\teil $16$ Hundertstel als Dezimalzahl? \leerfeld
\end{teile}
\end{aufgabe}

%% TODO Titel: Ich-kann-Satz fehlt in der Bank (Platzhalter: Kettenname) – Nr. 3
%% TODO Anweisung: ein Satz über der ersten Teilaufgabe fehlt in der Bank – Nr. 3
\begin{aufgabe}{Dezimalzahlen mit einer Zahl malnehmen und durch eine Zahl teilen, auch schriftlich oder mit dem Taschenrechner}
\begin{teile}
\teil $0{,}3 \cdot 4$? \leerfeld
\teil $1{,}25 \cdot 6$? \leerfeld
\end{teile}
\end{aufgabe}

%% TODO Titel: Ich-kann-Satz fehlt in der Bank (Platzhalter: Kettenname) – Nr. 4
%% TODO Anweisung: ein Satz über der ersten Teilaufgabe fehlt in der Bank – Nr. 4
\begin{aufgabe}{Dezimalzahlen vergleichen und ordnen (stellenweise, Nullen anhängen)}
\begin{teile}
\teil Größere Zahl: $0{,}7$ oder $0{,}9$? \leerfeld
\teil Ordne, mit der kleinsten beginnend: $0{,}45$; $0{,}405$; $0{,}54$ \leerfeld ; \leerfeld ; \leerfeld
\end{teile}
\end{aufgabe}

%% TODO Titel: Ich-kann-Satz fehlt in der Bank (Platzhalter: Kettenname) – Nr. 5
%% TODO Anweisung: ein Satz über der ersten Teilaufgabe fehlt in der Bank – Nr. 5
\begin{aufgabe}{Skala am Lineal und an der Uhr lesen (Untereinheiten)}
\begin{teile}
\teil Welche Zahl gehört zu $A$, welche zu $B$?

\zahlenstrahl[xmin=0,xmax=3,xstep=0.1,karo=0.5]{0.7/A, 2.3/B}
A: \leerfeld , B: \leerfeld
\teil Welche Zahl gehört zu $C$?

\zahlenstrahl[xmin=4,xmax=6,xstep=0.1,karo=0.6]{5.4/C}
C: \leerfeld
\end{teile}
\end{aufgabe}

%% TODO Titel: Ich-kann-Satz fehlt in der Bank (Platzhalter: Kettenname) – Nr. 6
%% TODO Anweisung: ein Satz über der ersten Teilaufgabe fehlt in der Bank – Nr. 6
\begin{aufgabe}{Längen im Koordinatensystem als Differenz von Koordinaten lesen (Abstand zweier Schnittstellen, Höhe zwischen Tiefpunkt und Scheitel, Funktionswert als Höhe), Flächeninhalte in Flächeneinheiten aus dem Integral oder der Figur und Volumina aus Querschnitt mal Länge}
\begin{teile}
\teil Punkte $P(2|1)$ und $Q(2|6)$ – senkrechter Abstand? \leerfeld
\teil Tiefpunkt $T(1|-1{,}5)$, Scheitel $S(0|2{,}25)$ – senkrechter Abstand? \leerfeld
\end{teile}
\end{aufgabe}

%% TODO Titel: Ich-kann-Satz fehlt in der Bank (Platzhalter: Kettenname) – Nr. 7
%% TODO Anweisung: ein Satz über der ersten Teilaufgabe fehlt in der Bank – Nr. 7
\begin{aufgabe}{Bruchteil einer Größe bilden (ein Viertel von einem Kilogramm; die Hälfte einer Stunde)}
\begin{teile}
\teil Ein Viertel von $8$ kg? \leerfeld[kg]
\teil Drei Viertel von $12$ l? \leerfeld[l]
\end{teile}
\end{aufgabe}

%% TODO Titel: Ich-kann-Satz fehlt in der Bank (Platzhalter: Kettenname) – Nr. 8
%% TODO Anweisung: ein Satz über der ersten Teilaufgabe fehlt in der Bank – Nr. 8
\begin{aufgabe}{Quadrat und dritte Potenz einer Länge mit Einheit (zwanzig Zentimeter zum Quadrat sind vierhundert Quadratzentimeter)}
\begin{teile}
\teil $3$ cm zum Quadrat – wie viel cm²? \leerfeld[cm²]
\teil $0{,}4$ m hoch drei – wie viel m³? \leerfeld[m³]
\end{teile}
\end{aufgabe}

%% TODO Titel: Ich-kann-Satz fehlt in der Bank (Platzhalter: Kettenname) – Nr. 9
%% TODO Anweisung: ein Satz über der ersten Teilaufgabe fehlt in der Bank – Nr. 9
\begin{aufgabe}{Prozentwert einer Größe berechnen (zehn Prozent eines Volumens)}
\begin{teile}
\teil $25\,\%$ von $36$ kg? \leerfeld[kg]
\teil $15\,\%$ von $2{,}4$ m³? \leerfeld[m³]
\end{teile}
\end{aufgabe}

%% TODO Titel: Ich-kann-Satz fehlt in der Bank (Platzhalter: Kettenname) – Nr. 10
%% TODO Anweisung: ein Satz über der ersten Teilaufgabe fehlt in der Bank – Nr. 10
\begin{aufgabe}{Formel nach einer Größe umstellen (aus einer Verhältnisgleichung die gesuchte Größe isolieren)}
\begin{teile}
\teil Löse: $x \cdot 4 = 28$ x = \leerfeld
\teil Löse: $1{,}6 \cdot x = 12$ x = \leerfeld
\end{teile}
\end{aufgabe}

%% TODO Titel: Ich-kann-Satz fehlt in der Bank (Platzhalter: Kettenname) – Nr. 11
%% TODO Anweisung: ein Satz über der ersten Teilaufgabe fehlt in der Bank – Nr. 11
\begin{aufgabe}{Fläche eines Rechtecks und eines Dreiecks berechnen (für den Materialbedarf)}
\begin{teile}
\teil Rechteck, $6$ m lang und $4$ m breit – Fläche? \leerfeld[m²]
\teil Rechteck, $2{,}4$ m lang und $1{,}5$ m breit – Fläche? \leerfeld[m²]
\end{teile}
\end{aufgabe}

%% TODO Titel: Ich-kann-Satz fehlt in der Bank (Platzhalter: Kettenname) – Nr. 12
%% TODO Anweisung: ein Satz über der ersten Teilaufgabe fehlt in der Bank – Nr. 12
\begin{aufgabe}{Multiplizieren und Dividieren mit zehn, hundert und tausend als Kommaverschiebung, auch bei Dezimalzahlen – Fehler finden}
\begin{teile}
\teil Lena rechnet: $8{,}4$ geteilt durch hundert ergibt $840$. Wo ist der Fehler?
\end{teile}
\end{aufgabe}

%% TODO Titel: Ich-kann-Satz fehlt in der Bank (Platzhalter: Kettenname) – Nr. 13
%% TODO Anweisung: ein Satz über der ersten Teilaufgabe fehlt in der Bank – Nr. 13
\begin{aufgabe}{Multiplizieren und Dividieren mit zehn, hundert und tausend als Kommaverschiebung, auch bei Dezimalzahlen – selbst rechnen}
\begin{teile}
\teil $6{,}5$ geteilt durch hundert? \leerfeld
\end{teile}
\end{aufgabe}
